\subsection{Compromisos y parámetros de diseño}

La arquitectura propuesta presenta dos ventajas principales. La primera es que permite reducir la huella de memoria sin reconstruir de forma
permanente el vector completo en su versión densa. La segunda es que introduce parámetros de control claros, como tamaño de bloque, capacidad de
caché, política de reemplazo y configuración del compresor, que permiten ajustar la relación entre ahorro espacial y costo temporal según el
patrón de acceso observado.

Estas mismas decisiones definen, sin embargo, un espacio de compromiso. Si los bloques son demasiado grandes, aumenta la compresibilidad
potencial, pero también el costo de recuperar datos para operaciones locales. Si los bloques son demasiado pequeños, disminuye el costo de acceso,
pero puede perderse oportunidad de compresión y aumentar la sobrecarga de metadatos. De forma similar, una caché muy reducida incrementa la tasa
de fallos y el número de descompresiones, mientras que una caché demasiado amplia amortigua accesos a costa de reservar más memoria densa activa.

La arquitectura también hereda una limitación más fundamental: su eficacia depende de la compresibilidad real del estado cuántico generado y del
patrón de compuertas que lo recorre. Cuando el estado presenta alta regularidad o redundancia, el diseño puede convertir esa estructura en ahorro
efectivo de memoria. En cambio, cuando el circuito genera amplitudes densas y poco compresibles o induce accesos dispersos sobre muchos bloques
distintos, la ventaja espacial puede disminuir y la sobrecarga temporal hacerse más visible.

En síntesis, la propuesta delimita una arquitectura de almacenamiento comprimido que no modifica el modelo de simulación de estado completo, pero
sí transforma la forma en que el vector de amplitudes reside y circula en memoria. Sobre esta base conceptual, el capítulo siguiente presenta la
materialización concreta de estas decisiones en el esquema de almacenamiento comprimido implementado sobre TMFQSfullstate con apoyo de Blosc2.
